% !TeX program = lualatex
% !TeX encoding = utf8
% !TeX spellcheck = uk_UA
% !BIB program = biber

\documentclass[biblatex]{ProblemBook}
\usepackage{lipsum}
\title{Електрика та магнетизм}
\def\subtitle{збірник задач}
\def\authors{С.~М.~Пономаренко}
%========================================================================================================
%
% !TeX program = lualatex
% !TeX encoding = utf8
% !TeX spellcheck = uk_UA
% !TeX root =../EMProblems.tex

%===============================================================================
%
%									      Палітурка
%
%===============================================================================

\newcommand{\CoverPage}{
	\begin{alwayssingle}
		\begin{center}
			\begin{flushright}\bfseries\sffamily
				\MakeUppercase{Міністерство освіти і науки України}\\
				КПІ ім. Ігоря Сікорського\\
			\end{flushright}
			\begin{tcolorbox}[titlepagestyle,
					toprule=0.10cm,
					bottomrule=0.10cm,
					overlay={%
						\node (picture) at ([xshift=4cm]frame.west) {\includegraphics[width=4cm]{logo_IPT}};
					}
			]%
			\begin{flushright}
				\large\bfseries\color{white}Фізико-технічний інститут
			\end{flushright}
			\end{tcolorbox}
			\vspace*{-2em}
		 	\begin{Large}\color{themecolordark!90!black}
			\begin{alignat*}{3}
				\vect{\nabla}&\times&&\Efield &&= -\dfrac{1}{c}\dfrac{\partial\Bfield}{\partial t} \\
				\vect{\nabla}&\,\cdot&&\Bfield &&= 0 \\
				\vect{\nabla}&\,\cdot&&\Dfield &&= 4\pi\rho \\
				\vect{\nabla}&\times&&\Hfield &&= \dfrac{4\pi}{c} \vect{j}+\dfrac{1}{c}\dfrac{\partial\Dfield}{\partial t}
			\end{alignat*}
			\end{Large}
			\vspace*{2em}
			\begin{tcolorbox}[
				titlepagestyle,
				toprule=0.15cm,
				bottomrule=0.15cm,
				top=1.3cm,
				bottom=0.7cm,
				overlay={%
				\node[%
							fill=white,
							rounded corners = 15pt,
							draw=themecolorlight,
							line width=0.15cm,
							inner sep=0pt,
							text width=8cm,
							minimum height=2cm,
							align=center,
							%anchor=east,
							font=\sffamily\bfseries\Large
						] (title) at (frame.north) {\authors};
				}
			]
			\centering
			\Huge\sffamily\bfseries\textcolor{white}{\realtitle}\\
			\huge\sffamily\bfseries\textcolor{white}{\subtitle}
			\end{tcolorbox}
			\vfill
		 	\begin{Large}\color{themecolordark!90!black}
			\begin{alignat*}{4}
					&\oiint\limits_S \Dfield\cdot d\vect{S} = 4\pi\iiint\limits_V\rho dV  \quad
					&&\oint\limits_L \Efield\cdot d\vect{r} &&&= - \frac1c \iint\limits_S \frac{\partial\Bfield}{\partial t}\cdot d\vect{S} \\
					&\oiint\limits_S \Bfield\cdot d\vect{S} = 0
					 \quad
					&&\oint\limits_L \Hfield\cdot d\vect{r} &&& =\dfrac{4\pi}{c} \iint\limits_S \vect{j}\cdot d\vect{S} \,+ \\
					& && &&& +\,\frac{1}{c} \iint\limits_S  \frac{\partial\Dfield}{\partial t}\cdot d\vect{S}
			\end{alignat*}
			\end{Large}
			\vfill
			\begin{tcolorbox}[titlepagestyle,
					toprule=0.10cm,
					bottomrule=0.10cm]
				\begin{center}\color{white}\bfseries\normalsize
					\MakeUppercase{Київ~2025} \\
%					КПІ ім. Ігоря Сікорського \\
%					2021
				\end{center}
			\end{tcolorbox}
		\end{center}
		\clearpage
	\end{alwayssingle}
\setcounter{page}{1}
}


%===============================================================================
%
%									      Титульна сторінка
%
%===============================================================================

\renewcommand\maketitle{
	\begin{alwayssingle}
		\begin{center}
				\MakeUppercase{Міністерство освіти і науки України}

				\bigskip
				\MakeUppercase{Національний технічний університет України}\\
				<<КИЇВСЬКИЙ ПОЛІТЕХНІЧНИЙ ІНСТИТУТ \\ імені ІГОРЯ СІКОРСЬКОГО>>
				\vspace*{100pt}

				{\Large\authors}
				\vspace*{50pt}

				{\Huge\sffamily\bfseries\realtitle}\\
				{\huge\sffamily\bfseries\subtitle}

			\vspace*{50pt}
			\begin{center}\itshape
				Рекомендовано Методичною радою КПІ ім. Ігоря Сікорського як навчальний посібник для здобувачів ступеня бакалавра за освітньою програмою <<Прикладна фізика>> спеціальності E6 <<Прикладна фізика та наноматеріали>>
			\end{center}

			\vfill
			\begin{center}
				\MakeUppercase{Київ} \\
				КПІ ім. Ігоря Сікорського \\
				2025
			\end{center}
		\end{center}
		\clearpage
	\end{alwayssingle}
}


%========================================================================================================
%
%									      Друга сторінка
%
%========================================================================================================
\newcommand\makeinfopage{
	\begin{alwayssingle}
		\noindent%
		\begin{minipage}[t]{0.5\textwidth}
				\begin{flushleft}
					УДК  537(075.8)\\
					ББК  22.31я73\\
					П 563
				\end{flushleft}
		\end{minipage}

%		\begin{minipage}[t]{0.55\textwidth}
%				\begin{flushleft}
%					Гриф надано Методичною радою КПІ ім. Ігоря Сікорського (протокол № 9/2019~від 30.05.2019~р.) за поданням Вченої ради Фізико-технічного інституту (протокол № 5/2019 від 24.04.2019 р.)
%				\end{flushleft}
%		\end{minipage}
		\bigskip\noindent%
		\begin{tabular}[t]{p{4.2cm}p{0.7\textwidth}}
	%				Автор та укладач:                        & \href{http://phes.ipt.kpi.ua/ponomarenko-sergij-mikolajovich}{\authors}, к.ф.-м.н., доцент кафедри  фізики енергетичних систем ФТІ \\
	%				                                         &                                                                                                                                   \\
			Рецензенти:                              & \href{http://imfn.lviv.ua/zf/?page_id=569}{І.~Р.~Зачек}, д.ф.-м.н., професор  кафедри загальної фізики Інституту прикладної математики та фундаментальних наук Національного університету <<Львівська політехніка>>                            \\
											 & \href{https://www.nas.gov.ua/UA/PersonalSite/Pages/default.aspx?PersonID=0000006279}{О.~М.~Кордюк}, д.ф.-м.н., член-кореспондент НАН України, директор Київського академічного університету                            \\
	                                     &                                                                                                                                   \\
			Відповідальний редактор: & С.~А.~Смирнов, к.ф.-м..н., доцент, голова методичної ради ФТІ
		\end{tabular}
		\vfill
		\begin{center}\itshape\small
				Гриф надано Методичною радою КПІ ім. Ігоря Сікорського (протокол №~9/2019~від 30.05.2019~р.) за поданням Вченої ради Фізико-технічного інституту (протокол №~5/2019 від 24.04.2019 р.)
		\end{center}
		\begin{center}
			\ifelectronic Електронне мережне видання \fi
			%\par {Версія від~\href{http://www.istpravda.com.ua/dates}{\today}} \par\else \par  \fi
		\end{center}
			\vfill
		\begin{center}
			\href{https://apd.ipt.kpi.ua/kafedra/stuff/ponomarenko}{\itshape Пономаренко Сергій Миколайович}, к.ф.-м.н., доцент
		\end{center}
			\vfill
		\begin{center}
			\LARGE\sffamily\realtitle \\
			\Large\sffamily\subtitle
		\end{center}
			\vfill
		\bigskip\noindent%
		\begin{flushleft}
			\begin{tabular}{lp{0.9\textwidth}}
				     & \textbf{\href{http://phes.ipt.kpi.ua/ponomarenko-sergij-mikolajovich}{\authors}}                                                                                                                                                                        \\
				П 563 & \hspace*{3ex} \realtitle : \subtitle{} [Електронний ресурс] : навчальний посібник / \authors{} --- К.:~КПІ ім. Ігоря Сікорського, 2025. --~\the\numexpr\getpagerefnumber{LastPage}-1\relax~с. -- Бібліогр.: с.~\pageref{BibPage}. \ifelectronic\relax\else-- 80~прим.\fi
			\end{tabular}
		\end{flushleft}
		\vfill

		Наведено \TotalValue{totalproblems} задачі з курсу <<Електрика та магнетизм>>. До більшості задач подано відповіді.

		Для студентів фізико-технічного інституту КПІ ім. Ігоря Сікорського, які навчаються за спеціальністю E6~<<Прикладна фізика та наноматеріали>>.

		\vfill

%		\begin{flushleft}\small
%			Ілюстративний матеріал підручника підготовлений за допомогою пакету \href{http://pgf.sourceforge.net}{TikZ/Pgf}. Верстка тексту проведена в видавничій системі \LaTeXe{} (компілятор Lua\LaTeX) на базі системи комп'ютерної верстки \TeX{} (Збірка  \href{https://www.tug.org/texlive/}{\TeX Live~2021}) з використанням оболонки \href{https://www.texstudio.org}{\TeX Studio}.
%		\end{flushleft}
%		\vfill

		\begin{flushleft}
			\begin{tabular}{p{\textwidth - 45ex}l}
				& \small\textcopyright\quad \authors, 2025~р. \\
				& \small \textcopyright\quad КПІ ім. Ігоря Сікорського (ФТІ), 2025~р.
			\end{tabular}
		\end{flushleft}
		\newpage%
	\end{alwayssingle}
}

%%%%%%%%%%%%%%%%%%%%%%%%%%%%%%%%%%%%%%%%%%%%%%%%%%%%%%%%%%%%%%%%%%%%%%%%%%%%
%%                                                                        %%
%%                              Last Page                                 %%
%%                                                                        %%
%%%%%%%%%%%%%%%%%%%%%%%%%%%%%%%%%%%%%%%%%%%%%%%%%%%%%%%%%%%%%%%%%%%%%%%%%%%%
\newcommand{\makelastpage}{%
\clearpage%
\thispagestyle{empty}%
\vspace*{0.4\textheight}
\begin{center}
	    \textbf{Пономаренко} Сергій Миколайович
\end{center}

\begin{center}\bfseries
    \Large\sffamily\realtitle \\
    \large\sffamily\subtitle
\end{center}

\begin{center}\itshape
    Комп'ютерне верстання в системі \LaTeXe{} С.\ М. Пономаренко
\end{center}

\vspace*{1em}
\begin{center}\small
Національний технічний університет України \\
<<Київський політехнічний інститут імені Ігоря Сікорського>> \\
Свідоцтво про державну реєстрацію: серія ДК № 5354 від 25.05.2017 р.\\
просп. Перемоги, 37, м. Київ, 03056
\end{center}

%\vfill
%
%\begin{center}
%    Видавництво <<Політехніка>> \kpishort \\
%    вул. Політехнічна, 14, корп. 15 \\
%    м. Київ, 03056 \\
%    Тел. (044) 204-81-78
%\end{center}
}






%					      Титульна сторінка
%
%========================================================================================================

\begin{document}
\CoverPage
\maketitle
\makeinfopage
\Opensolutionfile{answer}[Answers]

%
%
%
%\chapter{Chapter}
%
%%=========================================================
%\begin{problem}\label{prb:potential_6_2013-030-036_1}
%    У просторі між обкладками незаміщений плоский конденсатор вносять металеву пластину, що має заряд $Q$, так що між пластиною і обкладками конденсатора залишаються зазори $d_1$ та $d_2$ (рис.~\ref{potential_6_2013-030-036_1}). Площі пластини і обкладок конденсатора однакові і дорівнюють $S$. Визначити різницю потенціалів між обкладинками конденсатора. Крайовими ефектами нехтувати.
%\begin{solution}
%	$\Delta\phi = \frac{2\pi Q}{S} (d_2 - d_1)$.
%\end{solution}
%\end{problem}
%
%%=========================================================
%\begin{problem}\label{prb:potential_6_2013-030-036_4}
%    У плоскому конденсаторі на ліву обкладку поміщають заряд $+Q_1$, а на праву $+Q_2$. Всередину конденсатора паралельно обкладкам поміщають незаряджену металеву пластину (рис. 10). Які заряди будуть індуковані на лівій і правій поверхнях пластини, якщо $Q_2>Q_1$?
%\begin{solution}
%	$Q_{L} = \frac{Q_2 - Q_1}{2}$, $Q_{R} = -Q_{L}$.
%\end{solution}
%\end{problem}
%
%%=========================================================
%\begin{figure}[h!]\centering
%%---------------------------------------------------------
%\begin{minipage}[t]{0.45\linewidth}\centering
%	\begin{tikzpicture}
%		\draw[ultra thick] (0,0) -- +(0,2);
%		\draw[ultra thick] (3,0) -- +(0,2);
%		\draw[fill=gray!50, thick] (0.75,0) rectangle (1.75,2) ;
%		\node[above] at (1.5,2) {$+Q$};
%		\draw[latex-latex] (0,0.5) --node[above] {$d_1$} +(0.75,0);
%		\draw[latex-latex] (1.75,0.5) --node[above] {$d_2$} +(1.25,0);
%	\end{tikzpicture}
%\caption{До задачі~\ref{prb:potential_6_2013-030-036_1}}
%\label{potential_6_2013-030-036_1}
%\end{minipage}
%%---------------------------------------------------------
%\begin{minipage}[t]{0.45\linewidth}\centering
%	\begin{tikzpicture}
%		\draw[ultra thick] (0,0) -- +(0,2);
%		\draw[ultra thick] (3,0) -- +(0,2);
%		\draw[fill=gray!50, thick] (0.75,0) rectangle (1.75,2) ;
%		\node[above] at (0,2) {$+Q_1$};
%		\node[above] at (3,2) {$+Q_2$};
%		\draw[latex-latex] (0,0.5) --node[above] {$d_1$} +(0.75,0);
%		\draw[latex-latex] (1.75,0.5) --node[above] {$d_2$} +(1.25,0);
%	\end{tikzpicture}
%\caption{До задачі~\ref{prb:potential_6_2013-030-036_4}}
%\label{potential_6_2013-030-036_4}
%\end{minipage}
%%---------------------------------------------------------
%\end{figure}
%%=========================================================
%
%
%%=========================================================
%\begin{problem}\label{prb:potential_6_2013-030-036_2}
%    Три однакові нерухомі металеві пластини розташовані в повітрі на відстані $d_1$ і $d_2$ ($d_2> d_1$) одна від одної. Площа кожної з пластин дорівнює $S$ (рис.~\ref{potential_6_2013-030-036_2}). На середню пластину поміщають позитивний заряд $Q$. Пластини $1$ і $3$ не заряджені і підключені через ключ до резистору з невідомим опором, відмінним від нуля. Які заряди встановляться на пластинах $1$ і $3$ після довгого часу після замикання ключа? Крайовими ефектами нехтувати.
%\begin{solution}
%	$Q_1 = -\frac{Q}{2}\frac{d_2 - d_1}{d_2 + d_1}$, $Q_2 = - Q_1$.
%\end{solution}
%\end{problem}
%
%%=========================================================
%\begin{problem}\label{prb:potential_6_2013-030-036_5}
%    Три однакові металеві пластини розташовані в повітрі на рівних відстанях $d$ одна від одної. Площа кожної пластини дорівнює $S$ (рис.~\ref{potential_6_2013-030-036_5}). На пластині $1$ знаходиться позитивний заряд $Q$. Пластини $2$ і $3$ не заряджені і підключені через ключ до резистору з невідомим відмінним від нуля опором. Які набої встановляться на пластинах $2$ і $3$ після великого часу після замикання ключа?
%\begin{solution}
%	$Q_2 = -\frac{Q}{2}$, $Q_3 = \frac{Q}{2}$.
%\end{solution}
%\end{problem}
%
%%=========================================================
%\begin{figure}[h!]\centering
%%---------------------------------------------------------
%\begin{minipage}[t]{0.45\linewidth}\centering
%	\begin{tikzpicture}
%		\draw[ultra thick] (0,0) -- +(0,2) node[above] {$1$};
%		\draw[ultra thick] (1,0) -- node[pos=0.8, right] {$+Q$} +(0,2) node[above] {$2$};
%		\draw[ultra thick] (3,0) -- +(0,2) node[above] {$3$};
%		\draw[latex-latex] (0,0.5) --node[above] {$d_1$} +(1,0);
%		\draw[latex-latex] (1,0.5) --node[above] {$d_2$} +(2 ,0);
%		\draw (0,1) -- ++(-0.5,0) -- ++(0,-2) to[make contact] ++(1,0) to [resistor] ++(3,0) -- ++(0,2) -- ++(-0.5,0);
%	\end{tikzpicture}
%\caption{До задачі~\ref{prb:potential_6_2013-030-036_2}}
%\label{potential_6_2013-030-036_2}
%\end{minipage}
%%---------------------------------------------------------
%\begin{minipage}[t]{0.45\linewidth}\centering
%	\begin{tikzpicture}
%		\draw[ultra thick] (0,0) -- node[left] {$+Q$}+(0,2) node[above] {$1$};
%		\draw[ultra thick] (1,0) -- +(0,2) node[above] {$2$};
%		\draw[ultra thick] (2,0) -- +(0,2) node[above] {$3$};
%		\draw[latex-latex] (0,0.5) --node[above] {$d$} +(1,0);
%		\draw[latex-latex] (1,0.5) --node[above] {$d$} +(1 ,0);
%		\draw (1,0) -- ++(0,-0.25) -- ++(-1.5,0) -- ++(0,-0.75) to[make contact] ++(1,0) to [resistor] ++(2,0) -- ++(0,2) -- ++(-0.5,0);
%	\end{tikzpicture}
%\caption{До задачі~\ref{prb:potential_6_2013-030-036_5}}
%\label{potential_6_2013-030-036_5}
%\end{minipage}
%%---------------------------------------------------------
%\end{figure}
%%=========================================================
%
%%=========================================================
%\begin{problem}\label{prb:potential_6_2013-030-036_3}
%    Три тонкі незаряджені металеві пластини площею $S$ кожна розташовані на відстанях $d$ одна від одної, причому $d$ багато менше розмірів пластин. До пластин $2$ і $3$ під'єднали батарею з ЕРС $\EMF$ (рис.~\ref{potential_6_2013-030-036_3}). Пластині $1$ надали позитивний заряд $q_0$. Визначити заряд, який встановився на пластинах $2$ і $3$.
%\begin{solution}
%	$Q_2 = -\frac{q_0}{2} -\frac{\EMF S}{4\pi d}$, $Q_3 = \frac{q_0}{2} + \frac{\EMF S}{4\pi d}$.
%\end{solution}
%\end{problem}
%
%%=========================================================
%\begin{problem}\label{prb:potential_6_2013-030-036_6}
%    Три однакові нерухомі металеві пластини розташовані в повітрі на різних відстанях $d_1$ і $d_2$ ($d_2> d_1$) один від одного (рис.~\ref{potential_6_2013-030-036_6}). На середній пластині $2$ знаходиться позитивний заряд $Q$. спочатку не заряджені пластини $1$ і $3$ підключають через ключ до батареї з ЕРС $\EMF$ (рис. 12). Визначити встановилися заряди пластин $1$ і $3$ після замикання ключа.
%\begin{solution}
%	$Q_1 = -\frac{2\EMF S + 4\pi Q(d_2 - d_1)}{8\pi(d_1 + d_2)}$, $Q_3 = -Q_1$.
%\end{solution}
%\end{problem}
%%=========================================================
%\begin{figure}[h!]\centering
%%---------------------------------------------------------
%\begin{minipage}[t]{0.45\linewidth}\centering
%	\begin{tikzpicture}
%		\draw[ultra thick] (0,0) -- node[left] {$q_0$}+(0,2) node[above] {$1$};
%		\draw[ultra thick] (1,0) -- +(0,2) node[above] {$2$};
%		\draw[ultra thick] (2,0) -- +(0,2) node[above] {$3$};
%		\draw[latex-latex] (0,0.5) --node[above] {$d$} +(1,0);
%		\draw[latex-latex] (1,0.5) --node[above] {$d$} +(1 ,0);
%		\draw (1,0) -- ++(0,-1) to [battery={info'={$\EMF$},rotate=180}] ++(2,0) -- ++(0,2) -- ++(-1,0);
%	\end{tikzpicture}
%\caption{До задачі~\ref{prb:potential_6_2013-030-036_3}}
%\label{potential_6_2013-030-036_3}
%\end{minipage}
%%---------------------------------------------------------
%\begin{minipage}[t]{0.45\linewidth}\centering
%	\begin{tikzpicture}
%		\draw[ultra thick] (0,0) -- +(0,2) node[above] {$1$};
%		\draw[ultra thick] (1,0) -- node[pos=0.8, right] {$+Q$} +(0,2) node[above] {$2$};
%		\draw[ultra thick] (3,0) -- +(0,2) node[above] {$3$};
%		\draw[latex-latex] (0,0.5) --node[above] {$d_1$} +(1,0);
%		\draw[latex-latex] (1,0.5) --node[above] {$d_2$} +(2 ,0);
%		\draw (0,1) -- ++(-0.5,0) -- ++(0,-2) to[make contact] ++(1,0) to [battery={info'={$\EMF$}, rotate=180}] ++(3,0) -- ++(0,2) -- ++(-0.5,0);
%	\end{tikzpicture}
%\caption{До задачі~\ref{prb:potential_6_2013-030-036_6}}
%\label{potential_6_2013-030-036_6}
%\end{minipage}
%%---------------------------------------------------------
%\end{figure}
%%=========================================================

	\begin{center}
		\begin{tikzpicture}[scale=0.7]
			\begin{axis}[   axis y line = left,
					axis x line = bottom,
					grid = both,
					ylabel={$B$, кГс},
					xlabel={$H$, Е},
					xtick = {0,5,...,25},
					ytick = {0,5,...,25},
				]

				\addplot[thick, red, domain={0:25}, samples=100, name path = curve] {25*sqrt(1-((x-25)^2/25^2))};
				\addplot[blue, domain={0:20}, name path = line] {20 - x};
				\path [name intersections={of=line and curve, by=P}];
				\fill[red] (P) circle (0.05);
			\end{axis}
		\end{tikzpicture}
	\end{center}
%---------------------------------------------------------


\Closesolutionfile{answer}
\newpage
\input{Answers}
\end{document}






